\section{Numerical simulations}
\label{sec:rewrite-quantitative}

\subsection{Parameters and initial conditions}
\label{subsec:rewrite-rsi-activation}

In this section, I compare the transitional dynamics across four parameter configurations corresponding to the cases in Proposition~\ref{prop:rewrite-equilibrium-regimes}. I study an unanticipated activation of RSI at date zero. Before date zero, the economy already uses AI production services, but AI efficiency is fixed at $B_0$. For the chosen value of $B_0$, each economy starts on a labor-bottleneck balanced growth path. I compute this path with the research productivity parameter $\chi$ equal to zero to obtain the corresponding initial capital stock $K_0$ for each parameter configuration. Setting $\chi=0$ rules out improvements in AI efficiency, not the expansion of AI services through inference compute $U$. By treating activation as unanticipated, I exclude the advance adjustment of consumption and saving emphasized in the technology-news literature reviewed by \citet{beaudryportier2014}.

I compare four economies with $\sigma=0.9,1.0,1.1,1.5$, holding all other parameters and the initial values $A_0$, $N_0$, and $B_0$ fixed. For each economy, I choose $K_0$ to place it on its pre-RSI balanced growth path. In particular, for each elasticity, I choose the pre-event capital stock to satisfy
\begin{equation}
 r_0=\rho+\gamma,\qquad
 \frac{K_0}{Y_0}=\frac{\alpha}{\rho+\gamma+\delta}=3.30.
 \label{eq:rewrite-rsi-initial-ratio}
\end{equation}
This setup guarantees that $Y$, $X$, $U$, $w$, $r$, and $p_X$ behave continuously rather than jumping at date zero.

Table~\ref{tab:rewrite-rsi-parameters} reports the parameter values and initial stocks used for the simulations. I use standard values for $\alpha$, $\delta$, and $\rho$, while $n$ and $\gamma$ are guided by long-run population and productivity projections. The AI production weight $\omega_X=0.10$ and research exponent $\eta=0.20$ are illustrative choices, not empirical estimates. I choose them to satisfy $\omega_X+\eta<1$, the condition for the uncapped unit-elastic balanced-growth equilibrium in Proposition~\ref{prop:rewrite-uncapped-unit-bgp}; this is not a requirement for the simulations with an upper bound considered here. The common upper bound on AI efficiency satisfies $\mathcal B(1.5)<\overline B<\mathcal B(1.1)$. Thus, AI and labor are substitutes in both the $\sigma=1.1$ and $\sigma=1.5$ cases, but only the latter converges to the AI-dominated regime.

I study two levels of research productivity for each of the four elasticities of substitution. They both led to the same long-run growth rates in all variables, but differ in the speed at which growth accelerates once RSI becomes available. In the first one, higher-productivity scenario, $\chi$ is equal to 7.5. In the second one,  the lower-productivity scenario, $\chi$ is equal to 1.5. These are also illustrative choices, but their contrast show how different parameter values can lead to significant differences in the short-run impact of AI.

%In Equation~\eqref{eq:rewrite-rsi-initial-ratio}, $Y_0$ is obtained from the production and monopoly conditions with AI already present. Thus $K_0$ differs across elasticities, while the initial capital--output ratio and interest rate are common. With $A_0=N_0=1$ and the common $B_0$, the four capital stocks are approximately $3.434$, $4.649$, $5.037$, and $5.505$. These are calibrated alternative economies, not different outcomes of introducing AI into one economy with the same capital stock. The ratio is a theoretical BGP reference, not an exact empirical fit: the PWT 11.0 comparison yields a 2023 current-price capital--output ratio of about $3.61$ for the covered world sample and $3.42$ for the United States, with capital definitions that differ from the one-good model \citep{pwt110}.

% Edit Table 2 here: formatting, all rows, and notes are in this file.
% Display values only; this file does not set simulation inputs.
\begin{table}[H]
\centering
\caption{Simulation parameters and initial stocks}
\label{tab:rewrite-rsi-parameters}
\small\setstretch{1.0}
\setlength{\tabcolsep}{4pt}
\renewcommand{\arraystretch}{1.14}
\begin{tabular}{P{0.08\textwidth}P{0.12\textwidth}P{0.30\textwidth}P{0.36\textwidth}}
\toprule
Symbol & Value & Interpretation & Source\\
\midrule
$\alpha$ & 0.33 & Capital share & Standard calibration value.\\
$\delta$ & 0.05 & Depreciation rate & Standard calibration value.\\
$\rho$ & 0.04 & Discount rate & Standard calibration value.\\
$n$ & 0.003 & Population growth rate & Rounded constant-rate approximation to the UN 2024--2100 world projection \citep{unwpp2024}.\\
$\gamma$ & 0.01 & Growth of labor-augmenting technology & OECD labor-productivity projections \citep{oecdlongrun2025}.\\
$\eta$ & 0.20 & Research exponent & Illustrative.\\
%$(\omega_X,\omega_L)$ & (0.10, 0.90) & AI production and effective-labor weights & Illustrative.\\
$\omega_X$ & 0.10 & AI production weight & Illustrative.\\
$\omega_L$ & 0.90 & Effective-labor weight & $1-\omega_X$.\\
$\chi$ & 7.5 (higher);\newline 1.5 (lower) & Research productivity after activation & Illustrative.\\
$\sigma$ & 0.90;\newline 1.00;\newline 1.10;\newline 1.50. & Elasticity of substitution & Illustrative.\\
$\overline B$ & 1,361 & AI-efficiency upper bound & $1.10 \times \mathcal B(1.50)$.\\
\midrule
%$A_0,N_0$ & 1; 1 & Initial technology and population & Normalizations.\\
$A_0$ & 1 & Initial labor-efficiency units & Normalization.\\
$N_0$ & 1 & Initial population & Normalization.\\
$B_0$ & 13.6 & Initial AI efficiency & $0.01 \times \overline B$.\\
$K_0$ & 3.43;\newline 4.65;\newline  5.04;\newline 5.50. & Initial capital & Adjusted to match $K_0/Y_0 =3.30$, the pre-RSI balanced growth path ratio.\\
%$K_0/Y_0$ & 3.30 & Initial capital--output ratio & Derived from the pre-RSI balanced growth path.\\
\bottomrule
\end{tabular}
\par\vspace{4pt}
%\parbox{0.98\textwidth}{\footnotesize\emph{Notes:} %Rates are annual. Values are rounded for display; the simulations retain their original precision. $L_0=N_0$. %Initial consumption $C_0$ and the developer's shadow value $q_0$ are selected endogenously by the BVP for each path, not supplied as inputs. $K_0/Y_0$ is a derived ratio, measured in years of date-zero output, not an additional parameter. Lists of four ratios follow the displayed order of $\sigma$.}
\end{table}



%Before activation, consumption satisfies
%\begin{equation}
% C(0^-)=Y_0-U_0-(\delta+n+\gamma)K_0>0.
% \label{eq:rewrite-rsi-pre-consumption}
%\end{equation}
%Only $K_0$ and $B_0$, together with the exogenous $A_0,N_0$, are inherited by the post-event equilibrium. The BVP selects $C(0^+)$ and $q(0^+)$ anew. 
%The pre-event BGP requires an unanticipated activation: announcing RSI earlier would generally change saving and consumption before date zero, as in the technology-news mechanisms reviewed by \citet{beaudryportier2014}.

%Unlike the scenario in \citet{jonestonetti2026}, where endogenous automation is already underway, this exercise studies an unanticipated activation of RSI and so subsequent AI-efficiency improvements are endogenous. However, the exercise omits the costs of expanding compute capacity as in \citep{erdiletal2025gate} and the complementary intangible investment studied by \citet{brynjolfssonrocksyverson2021}. Those mechanisms can delay adjustment or measured productivity gains. 

I do not model a separate stock of computing hardware or the costs of
rapidly expanding it. \citet{erdiletal2025gate} incorporate adjustment costs
for compute capacity and other physical capital. Introducing such frictions
could change the initial reallocation of resources and the timing of the
transition shown here.

Appendix~\ref{app:rewrite-algorithm}
describes the algorithm used to find the equilibrium and presents the links to replication files.
